\section{An incidence criterion and its embedded local groups}
\label{sec:criterion}

We now leave affine geometry. The following criterion is the algebraic part of the construction. Its
input consists of a finite incidence system and vectors satisfying two
matrix identities. All groups and group-algebra elements in its conclusion
are constructed in this section and Section~\ref{sec:packing}; no further geometric assumption enters those arguments.

\begin{theorem}[Incidence criterion]
\label{thm:criterion}
Let $\ell$ be an odd prime, let $K$ be a finite field of characteristic two
containing a primitive $\ell$th root of unity $\zeta$, and put
$\kappa=\ell^2$. Let $V$ be a finite set of cardinality $t$, and let
$\mathcal F$ be a finite family of nonempty subsets of $V$. Assume that
every $v\in V$ belongs to exactly $\ell+1$ members of $\mathcal F$ and
that the bipartite incidence graph of $V$ and $\mathcal F$ has girth at
least $12$.

Suppose that an integer $m$ and column vectors $x_v\in K^m$, $v\in V$,
satisfy
\begin{equation}
\label{eq:criterion-size}
0<m<t\le m\kappa
\end{equation}
and
\begin{equation}
\label{eq:criterion-pairing}
\sum_{v\in V}x_vx_v^{\mathsf T}=-I_m,
\qquad
\sum_{v\in f}x_vx_v^{\mathsf T}=0
\quad\text{for every }f\in\mathcal F.
\end{equation}
Then there are a finitely presented group $G$ containing an element of order $\ell$ and finite sums
$a_{\mathrm{out}},b_{\mathrm{out}}\in K[G]$ such that
\[
a_{\mathrm{out}}b_{\mathrm{out}}=1,
\qquad
b_{\mathrm{out}}a_{\mathrm{out}}\ne1.
\]
They are given by the constructions and formulas below.
\end{theorem}

The hypotheses have separate roles. Incidence girth preserves the local
subgroups, while Equation~\eqref{eq:criterion-pairing} produces a
rectangular matrix identity. The upper bound $t\le m\kappa$ permits
character packing into square matrices; the strict inequality $m<t$
then detects their reverse-product defect in an embedded subgroup algebra.
An explicit finite group supplies matrix units that turn these matrices
into the promised elements.

\subsection{The glued group}

We first realize the incidence data by finite subgroups of one group.
For each $v\in V$, let $K_v^0$ be a coordinatized copy of the additive
group of $\F_\ell^2$, and for each $f\in\mathcal F$, let $L_f^0$ be a
copy of the additive group of $\F_\ell$.  Since exactly $\ell+1$
hyperedges contain $v$, choose injections
\[
  \iota_{vf}:L_f^0\longrightarrow K_v^0\qquad(v\in f)
\]
whose images are the $\ell+1$ distinct one-dimensional subspaces of
$K_v^0$.  These are homomorphisms of additive groups.  In particular,
their images partition $K_v^0\setminus\{0\}$ outside their common
zero element.

Introduce a generator $s_{f,c}$ for each $f\in\mathcal F$ and each
$0\ne c\in L_f^0$.  A signed word
\[
  s_{f_1,c_1}^{\varepsilon_1}\cdots
  s_{f_k,c_k}^{\varepsilon_k},\qquad \varepsilon_i\in\{1,-1\},
\]
is a defining relator whenever some $v\in V$ belongs to every $f_i$
and
\[
  \sum_{i=1}^k\varepsilon_i\iota_{vf_i}(c_i)=0
  \quad\text{in }K_v^0.
\]
Let $N$ be the group with these generators and relators.  It is
finitely generated, since the incidence data and all the local groups
are finite.

There is a homomorphism $\theta_v:K_v^0\to N$ determined by
\[
  \theta_v(0)=1,\qquad
  \theta_v\bigl(\iota_{vf}(c)\bigr)=s_{f,c}\quad(c\ne0).
\]
Indeed, each nonzero element has a unique expression on the left.
For $a,b\in K_v^0$, the word expressing
$\theta_v(a)\theta_v(b)\theta_v(a+b)^{-1}$ is a defining relator,
after omitting factors corresponding to zero elements.

The same group is presented by the finite sublist of defining relators
of length at most three. These include every local multiplication-table
relation just displayed, as well as the local inverse relations. Within
one $K_v^0$, successive applications of those relations reduce any
signed word to the generator for its sum, or to the empty word when
the sum is zero. Thus they imply every defining relator of $N$, and
$N$ is finitely presented.
We retain the all-word formulation because it makes the planar proof
below more transparent.

\begin{lemma}\label{lem:embedding}
If the incidence graph of $(V,\mathcal F)$ has girth at least $12$,
then every homomorphism $\theta_v:K_v^0\to N$ is injective.
\end{lemma}

\begin{proof}
Every nonzero element of $K_v^0$ maps to a single generator $s_{f,c}$.
It therefore suffices to prove that no such generator is the identity
in $N$.  Suppose, to the contrary, that $s_{f,c}=1$, with $c\ne0$.
We will encode this equality by a planar graph with zero local sums
at every vertex except one.  Reducing the graph will make the local
group relations force lower bounds on vertex degrees, while incidence
girth forces lower bounds on facial lengths.  These bounds will
contradict Euler's formula even with the one exceptional vertex.

\emph{From a word equality to a planar graph.}
The one-letter word $s_{f,c}$ can be changed to the empty word by
finitely many insertions and deletions of contiguous defining
relators.  To justify this form of reduction, take the congruence on
signed words generated by those operations.  Every formal inverse
pair is a defining relator, since its hyperedge has an incident
vertex.  The quotient is consequently a group, and its universal
property is exactly that of the presentation of $N$.  Thus equality
in $N$ is equality under this congruence.  Equivalently, each factor
$u r^{\pm1}u^{-1}$ in a finite normal-closure expression can be
inserted by first forming $u u^{-1}$ through nested inverse-pair
insertions and then inserting $r^{\pm1}$ between its two halves.
Inverses of defining relators are again defining relators.

Draw the successive words on horizontal levels, treating each signed
letter as an individual token.  Between two levels, connect the
tokens that persist in their unchanged order.  An inserted or deleted
block begins or ends its tokens at a new black vertex, labelled by a
point $v$ at which the corresponding relator has zero sum.  This can
be drawn without crossings: the affected block is contiguous, so its
star occupies a single gap between the other order-preserving wires.
Joining the finitely many strips gives a planar drawing.

Every inserted token $s_{i,y}^{\varepsilon}$ has one creation vertex
and one deletion vertex.  Subdivide its full lifetime wire by a white
vertex labelled $i$, and give its two edges weights $\varepsilon y$
at the creation end and $-\varepsilon y$ at the deletion end.  The
initial token instead has an exceptional white endpoint labelled $f$,
joined to its deletion vertex by an edge of weight $-c$.  We obtain
a finite planar bipartite multigraph with black labels in $V$ and
white labels in $\mathcal F$.  An edge joining labels $v$ and $i$
has a nonzero weight in $L_i^0$, and $v\in i$.  Its balances are
\begin{equation}\label{eq:embedding-balances}
\begin{aligned}
  &\sum_{a\ni x}\iota_{v,i(a)}\bigl(\omega(a)\bigr)=0
    &&\text{at each black vertex $x$ labelled $v$},\\
  &\sum_{a\ni y}\omega(a)=0
    &&\text{at each ordinary white vertex $y$},
\end{aligned}
\end{equation}
where $\omega(a)$ is the weight of edge $a$ and $i(a)$ is its white
label.  The exceptional white vertex has nonzero sum.  At creation
vertices the first balance is the relator sum, and at deletion
vertices it is its negative.  Ordinary white sums vanish because
their two weights are opposite.  Retain only the connected component
of the exception.

\emph{Reducing the planar graph.}
Among all connected planar bipartite multigraphs with the indicated
labels, nonzero edge weights, and balances
\eqref{eq:embedding-balances}, and with exactly one exceptional white
vertex of nonzero sum, choose one minimizing lexicographically the
number of vertices and then the number of edges.  We do not require
this comparison graph to come from a word reduction.  Its exceptional
vertex has positive degree, so the graph has an edge.

There are no parallel edges.  Two parallel edges have weights in the
same $L_i^0$ and the same map $\iota_{vi}$ at their black endpoint.
Replace them by one edge carrying their sum, or delete both if that
sum is zero, and again retain the component of the exception.
Additivity preserves every balance and preserves the exceptional
nonzero sum.  This gives a smaller comparison graph, a contradiction.

Fix a spherical embedding of the minimal graph.  We claim that, at
each vertex of degree at least two, cyclically consecutive neighbors
have distinct labels.  If two such neighbors have the same label,
they are distinct vertices because parallel edges have been removed.
Join them by an auxiliary arc following one side of the two incident
edges and the empty sector between those edges, and contract the
arc.  A sufficiently narrow neighborhood of this two-edge path
makes the arc disjoint from the rest of the graph.  This construction
is local and remains valid when either edge is a bridge.

The contraction merges vertices of the same color and label.  Each
original edge keeps its weight and its opposite-color endpoint; hence
the resulting graph is still bipartite, and every edge weight is
nonzero.  Parallel edges created by the contraction are allowed in
the comparison class.  When black vertices are merged, their
balances add in the same $K_v^0$, while white balances are unchanged.
When white vertices are merged, their balances add in the same
$L_i^0$, while every black incidence map and weighted sum is
unchanged.  If the exception is merged, its sum is increased by zero.
Thus the contraction gives a comparison graph with fewer vertices,
proving the claim.

\begin{figure}[htbp]
  \centering
  \begin{tikzpicture}[
    x=1cm, y=1cm, font=\small,
    merge white/.style={circle, draw, fill=white, line width=.6pt,
      minimum size=5.5pt, inner sep=0pt},
    merge black/.style={circle, draw, fill=black,
      minimum size=4.5pt, inner sep=0pt},
    merge edge/.style={line width=.6pt},
    merge auxiliary/.style={densely dashed, line width=.55pt}
  ]
    \begin{scope}
      \coordinate (merge-before-p-position) at (-1.1,1.65);
      \coordinate (merge-before-q-position) at (1.1,1.65);
      \coordinate (merge-before-x-position) at (0,0);
      \draw[merge edge]
        (-1.55,2.2) -- (merge-before-p-position)
        -- (merge-before-x-position) -- (merge-before-q-position)
        -- (1.55,2.2);
      \draw[merge auxiliary]
        (merge-before-p-position)
        .. controls (-.92,1.47) and (-.15,.43) .. (0,.43)
        .. controls (.15,.43) and (.92,1.47) ..
        (merge-before-q-position);
      \node[merge white, label=above:{$f$}, label=left:{$p$}]
        (merge-before-p) at (merge-before-p-position) {};
      \node[merge white, label=above:{$f$}, label=right:{$q$}]
        (merge-before-q) at (merge-before-q-position) {};
      \node[merge black, label=below:{$v$}, label=left:{$x$}]
        (merge-before-x) at (merge-before-x-position) {};
      \node at (-.76,.78) {$a$};
      \node at (.76,.78) {$b$};
    \end{scope}

    \draw[-{Stealth[length=5pt]}, line width=.6pt]
      (1.95,1) -- node[above=3pt] {contract} (3.35,1);

    \begin{scope}[xshift=5.1cm]
      \coordinate (merge-after-r-position) at (0,1.65);
      \coordinate (merge-after-x-position) at (0,0);
      \draw[merge edge]
        (-.75,2.2) -- (merge-after-r-position) -- (.75,2.2);
      \draw[merge edge]
        (merge-after-x-position)
        .. controls (-.63,.45) and (-.63,1.2) ..
        (merge-after-r-position);
      \draw[merge edge]
        (merge-after-x-position)
        .. controls (.63,.45) and (.63,1.2) ..
        (merge-after-r-position);
      \node[merge white, label=above:{$f$}, label=right:{$r$}]
        (merge-after-r) at (merge-after-r-position) {};
      \node[merge black, label=below:{$v$}, label=left:{$x$}]
        (merge-after-x) at (merge-after-x-position) {};
      \node at (-.72,.82) {$a$};
      \node at (.72,.82) {$b$};
    \end{scope}
  \end{tikzpicture}
  \caption{The white vertices $p,q$, both labelled $f$, merge to $r$
    along the dashed auxiliary arc in the empty sector between their
    consecutive edges at the black vertex $x$, labelled $v$.
    The edges of weights $a,b$ remain distinct parallel edges.}
  \label{fig:merge}
\end{figure}


\emph{Degrees, faces, and Euler's formula.}
An ordinary white vertex cannot have degree one, since its incident
weight would be nonzero.  A black vertex cannot have degree one
either, since the maps $\iota_{vi}$ are injective.  If a black vertex
had degree two, its two white neighbors would have distinct labels
by the preceding claim.  Its zero balance would then equate two
nonzero vectors in distinct one-dimensional subspaces of $K_v^0$,
which is impossible.  Consequently every black degree is at least
three, and every white degree is at least two except possibly the
exceptional degree, which is at least one.

Consider the facial boundary walks of this connected plane graph,
counting edge sides with multiplicity.  At a turn through a vertex
of degree at least two, the entering and leaving edges are distinct
and consecutive in its rotation.  Their opposite endpoints have
different labels.  Thus the image of the boundary walk in the
incidence graph has no immediate backtracking at that turn.  This
also applies to the cyclic closing turn of the walk.  Bridges cause
no difficulty: their two sides are counted separately, and a facial
walk reverses immediately along an edge only at a degree-one vertex.
There is at most one such vertex, the exceptional white vertex, and
it belongs to just one face.

It follows that all but at most one face have boundary walks whose
images are nonempty closed walks without immediate backtracking in
the incidence graph.  Such a walk contains a cycle, so each of these
faces has length at least $12$.  Write $n_{\rm b},n_{\rm w},n_{\rm f}$
for the numbers of black vertices, white vertices, and faces, and
$E$ for the number of edges.  The degree bounds and the total of
$2E$ edge sides give
\[
  n_{\rm b}\le\frac E3,\qquad
  n_{\rm w}\le\frac{E+1}{2},\qquad
  n_{\rm f}\le\frac E6+1.
\]
Therefore
\[
  n_{\rm b}+n_{\rm w}+n_{\rm f}
  \le E+\frac32<E+2,
\]
contrary to Euler's formula for a connected plane graph.  The
assumed word equality is impossible, and every $\theta_v$ is
injective.
\end{proof}

We henceforth identify $K_v^0$ with its image $K_v\le N$, and identify
$L_f^0$ with the subgroup $L_f\le N$ given by $c\mapsto s_{f,c}$,
with $0$ mapped to the identity.  The latter is injective by choosing
any point incident to $f$ and applying Lemma~\ref{lem:embedding}.
The shared generators identify this same subgroup inside every
$K_v$ with $v\in f$.  In particular, $|K_v|=\ell^2$, $|L_f|=\ell$,
and the incident subgroups $L_f$ partition $K_v$ outside the
identity.  No assertion about intersections between other local
subgroups is needed.


\section{From subgroup averages to the scalar defect}
\label{sec:packing}

The embedded groups now allow legitimate subgroup averages. We first obtain a rectangular identity, then pack its intermediate coordinates into square matrices, and finally encode those matrices as scalar elements of a group algebra. The strict size inequality will detect the reverse defect after packing.

\subsection{Subgroup averages and a rectangular identity}

We prove Theorem~\ref{thm:criterion} using the group $N$ and subgroups
$K_v,L_f\le N$ supplied by Lemma~\ref{lem:embedding}. We retain the
coordinates on $K_v$ transported from $K_v^0=\F_\ell^2$. Fix an
ordering of $V$, and use it for all matrix indices indexed by $V$.
For $v\in V$ put
\begin{equation}
\label{eq:local-average}
P_v=\kappa^{-1}\sum_{u\in K_v}u\in K[N].
\end{equation}
The integer $\kappa$ is invertible in $K$, since $\ell$ is odd. For each
$w\in K_v$ there are exactly $\kappa$ pairs $(u,u')\in K_v^2$ with
$uu'=w$, so $P_v^2=P_v$.

At a fixed $v$, the subgroups $L_f$ for $f\ni v$ run through all the
one-dimensional subspaces of the additive group $K_v$. Their nonidentity
elements partition $K_v\setminus\{1\}$. Consequently
\begin{equation}
\label{eq:subgroup-decomposition}
\sum_{u\in K_v}u
=1+\sum_{f\ni v}\left(\sum_{u\in L_f}u-1\right).
\end{equation}
This equality is internal to $K_v$; no assertion about intersections
between different $K_v$ is needed.

Write $W_v=x_vx_v^{\mathsf T}$. Define matrices over $K[N]$ by
\begin{equation}
\label{eq:rectangular-definitions}
U\in\Mat_{m\times t}(K[N]),\qquad
T\in\Mat_{t\times m}(K[N]),\qquad
U_{-,v}=-\kappa x_vP_v,\quad
T_{v,-}=P_vx_v^{\mathsf T},
\end{equation}
and set $P=\diag_{v\in V}(P_v)$. Since the entries of each $x_v$
belong to the coefficient field, they commute with all group-algebra
elements. Idempotence gives $UP=U$ and $PT=T$. Moreover,
\begin{align*}
UT
&=-\kappa\sum_{v\in V}W_vP_v^2
 =-\sum_{v\in V}W_v\sum_{u\in K_v}u\\
&=-\sum_{v\in V}W_v
 -\sum_{f\in\mathcal F}
   \left(\sum_{v\in f}W_v\right)
   \left(\sum_{u\in L_f}u-1\right)
 =I_m,
\end{align*}
where we used Equations~\eqref{eq:subgroup-decomposition} and
\eqref{eq:criterion-pairing}. We have proved
\begin{equation}
\label{eq:rectangular-identities}
UP=U,\qquad PT=T,\qquad UT=I_m.
\end{equation}

\subsection{A group extension implementing character twists}

Equation~\eqref{eq:rectangular-identities} gives the forward identity,
but the factors are rectangular. We next arrange enough orthogonal
character idempotents to place their $t$ intermediate coordinates into
$m$ slots. The extension will implement the necessary changes of character
by conjugation while preserving the original subgroup algebra.

Let $z$ generate a cyclic group of order $\ell$, and put
$H=N\times\langle z\rangle$. Choose $v_*\in V$ and put $K_*=K_{v_*}$.
For every $v\in V$ let $\phi_v:K_v\longrightarrow K_*$ be the
isomorphism taking each element to the element with identical
$\F_\ell^2$ coordinates. There are exactly $\kappa$ linear forms
$K_*\longrightarrow\F_\ell$. By Equation~\eqref{eq:criterion-size},
we may assign to each $v\in V$ a slot $s(v)\in\{1,\ldots,m\}$ and a
linear form $\alpha_v:K_*\longrightarrow\F_\ell$ so that the pairs
$(s(v),\alpha_v)$ are distinct. In particular, the forms assigned within
any one slot are distinct.

We use a multiple HNN extension~\cite[Theorem~II]{HNN1949}, verifying
the required base-group embedding directly. Adjoin to $H$ generators
$\tau_v$, $v\in V$, with relations
\begin{equation}
\label{eq:twisted-conjugation}
\tau_v(uz^r)\tau_v^{-1}
 =\phi_v(u)z^{r+\alpha_v(\phi_v(u))}
\qquad(u\in K_v,\ r\in\mathbb Z/\ell).
\end{equation}
Denote the resulting group by $L$. Exponents in $\F_\ell$ are
interpreted in $\mathbb Z/\ell$. For each $v$, the specified map
\[
\psi_v:K_v\times\langle z\rangle\longrightarrow
K_*\times\langle z\rangle,
\qquad
uz^r\longmapsto\phi_v(u)z^{r+\alpha_v(\phi_v(u))},
\]
is a group isomorphism. Its inverse takes $wz^r$ to
$\phi_v^{-1}(w)z^{r-\alpha_v(w)}$.

We verify directly that $H$ embeds in $L$. Let $H$ act on itself by left
multiplication. Restricted to either the domain or the image of
$\psi_v$, this is a disjoint union of regular actions of a subgroup of
order $\ell^3$. These two collections of orbits have the same
cardinality. Indeed, if $H$ is finite, both have $|H|/\ell^3$ orbits;
if $H$ is infinite, both have countably infinitely many, since $H$ is
countable and the subgroups are finite. Pair these orbits and choose
representatives $h_i,k_i$ so that they are respectively
$(K_v\times\langle z\rangle)h_i$ and
$(K_*\times\langle z\rangle)k_i$. The bijection
\[
F_v(ah_i)=\psi_v(a)k_i
\qquad(a\in K_v\times\langle z\rangle)
\]
of the set $H$ conjugates left multiplication by $a$ to left
multiplication by $\psi_v(a)$. Assign $F_v$ to $\tau_v$ for each $v$.
These assignments, together with the left regular action of $H$, satisfy
all the relations defining $L$: no defining relation couples two distinct
stable letters, so overlapping associated subgroups impose no additional
compatibility condition on the independently chosen $F_v$. They therefore give a permutation
representation of $L$ whose restriction to $H$ is faithful. This proves
the required embedding. In particular, $K[H]$ embeds in $K[L]$.

Taking $u=1$ in Equation~\eqref{eq:twisted-conjugation} shows that every
$\tau_v$ centralizes $z$. Since $z$ is already central in $H$, it is
central in $L$. Define
\begin{equation}
\label{eq:central-character-idempotent}
e=\ell^{-1}\sum_{r=0}^{\ell-1}\zeta^{-r}z^r\in K[L].
\end{equation}
This is a nonzero central idempotent and satisfies $ez=\zeta e$.
For completeness, in the convolution square each coefficient of $z^a$
is $\ell^{-2}\ell\zeta^{-a}$, which proves $e^2=e$; reindexing the
sum proves $ez=\zeta e$. The identity coefficient is $\ell^{-1}\ne0$
because $\langle z\rangle$ remains embedded.

For a linear form $\alpha:K_*\longrightarrow\F_\ell$, put
\begin{equation}
\label{eq:local-character-idempotent}
Q_\alpha=\kappa^{-1}\sum_{u\in K_*}\zeta^{\alpha(u)}u.
\end{equation}
Using additive notation for the coordinates of $K_*$, the coefficient
of $w$ in $Q_\alpha Q_\beta$ is
\[
\kappa^{-2}\zeta^{\beta(w)}
\sum_{u\in K_*}\zeta^{\alpha(u)-\beta(u)}.
\]
If $\alpha=\beta$, the inner sum is $\kappa$. Otherwise choose $u_0$
with $\zeta^{\alpha(u_0)-\beta(u_0)}\ne1$. Translating the summation
variable by $u_0$ multiplies the sum by this scalar while leaving it
unchanged, so the sum is zero. Thus
\begin{equation}
\label{eq:character-products}
Q_\alpha Q_\beta=
\begin{cases}
Q_\alpha,&\alpha=\beta,\\
0,&\alpha\ne\beta.
\end{cases}
\end{equation}
The conjugation relations and $ez^r=\zeta^re$ now give
\begin{equation}
\label{eq:twisted-projection}
\tau_v eP_v\tau_v^{-1}
=\kappa^{-1}\sum_{u\in K_v}
 e\phi_v(u)z^{\alpha_v(\phi_v(u))}
=eQ_{\alpha_v}.
\end{equation}

\subsection{Packing into square matrices}

The orthogonality in Equation~\eqref{eq:character-products} now lets
different coordinates share a slot without cross terms. The following
matrices implement the assignment of each coordinate $v$ to its slot
$s(v)$.

Define $C\in\Mat_{m\times t}(K[L])$ and
$J_0\in\Mat_{t\times m}(K[L])$ by
\begin{equation}
\label{eq:packing-matrices}
C_{s,v}=
\begin{cases}
\tau_veP_v,&s=s(v),\\
0,&s\ne s(v),
\end{cases}
\qquad
(J_0)_{v,s}=
\begin{cases}
eP_v\tau_v^{-1},&s=s(v),\\
0,&s\ne s(v).
\end{cases}
\end{equation}
We claim that
\begin{equation}
\label{eq:packing-identity}
J_0C=eP.
\end{equation}
The diagonal entry indexed by $v$ is $eP_v$. If $v\ne w$ have different
slots, the corresponding entry is zero. If they have the same slot,
Equation~\eqref{eq:twisted-projection} rewrites that entry as
\[
eP_v\tau_v^{-1}\tau_weP_w
=\tau_v^{-1}eQ_{\alpha_v}eQ_{\alpha_w}\tau_w=0,
\]
because $\alpha_v\ne\alpha_w$. This proves the claim.

Set
\begin{equation}
\label{eq:square-elements}
a_{\mathrm{mat}}=UJ_0+(1-e)I_m,
\qquad
b_{\mathrm{mat}}=CT+(1-e)I_m.
\end{equation}
Every entry of $C$ and $J_0$ contains the central factor $e$.
Consequently the cross terms with $1-e$ vanish, and
Equations~\eqref{eq:rectangular-identities} and
\eqref{eq:packing-identity} imply
\begin{equation}
\label{eq:square-forward-product}
a_{\mathrm{mat}}b_{\mathrm{mat}}
=U(eP)T+(1-e)I_m
=eI_m+(1-e)I_m=I_m.
\end{equation}

\subsection{Detecting the reverse defect}

Put $\Delta=b_{\mathrm{mat}}a_{\mathrm{mat}}-I_m$. We detect its
nonvanishing by multiplying on the left by $J_0$ and on the right by
$C$. This brings the calculation back from $K[L]$ to its embedded
subgroup algebra $K[H]$, where a scalar specialization is available.
Direct multiplication
gives $\Delta=CTUJ_0-eI_m$, and hence
\begin{equation}
\label{eq:reverse-sandwich}
J_0\Delta C=(eP)TU(eP)-eP=e(TU-P).
\end{equation}
In the last equality the order of multiplication is essential:
$(eP)TU(eP)=e(PT)(UP)=eTU$. The right-hand side of
Equation~\eqref{eq:reverse-sandwich} already belongs to
$\Mat_t(K[H])$.

There is a $K$-algebra homomorphism
\begin{equation}
\label{eq:subgroup-specialization}
\varepsilon:K[H]\longrightarrow K,
\qquad
\varepsilon(nz^r)=\zeta^r\quad(n\in N).
\end{equation}
It is well-defined because $H=N\times\langle z\rangle$.
It sends $e$ to $1$ and each $P_v$ to $1$.
Let $X\in\Mat_{m\times t}(K)$ be the matrix with columns $x_v$.
Since the characteristic is two and $\kappa$ is odd,
\[
\varepsilon(U)=X,\qquad
\varepsilon(T)=X^{\mathsf T},\qquad
XX^{\mathsf T}=I_m.
\]
Thus
\begin{equation}
\label{eq:specialized-defect}
\varepsilon\bigl(e(TU-P)\bigr)=X^{\mathsf T}X-I_t.
\end{equation}
Because $t>m$, choose a nonzero $y\in\ker X$. Then
\begin{equation}
\label{eq:nonzero-witness}
\varepsilon\bigl(e(TU-P)\bigr)y=-y\ne0.
\end{equation}
One explicit choice is the first nonzero column of $I_t-X^{\mathsf T}X$:
this matrix is nonzero because $\rank(X^{\mathsf T}X)\le m<t$, and
$X(I_t-X^{\mathsf T}X)=0$. This choice uses only finite-field
arithmetic on the given vectors.

Equation~\eqref{eq:nonzero-witness} proves that $e(TU-P)$ is nonzero in
$\Mat_t(K[H])$. The embedding $K[H]\subseteq K[L]$ and
Equation~\eqref{eq:reverse-sandwich} therefore show that $\Delta\ne0$.
The homomorphism $\varepsilon$ has been applied only to elements of
$K[H]$; no extension of it to $K[L]$ is used.

\subsection{From matrices to two group-algebra elements}

The square factors give a failure of stable finiteness for $K[L]$.
To obtain scalar factors, we embed $\Mat_m(K[L])$ in an idempotent
corner of $K[L\times F]$ for a suitable finite group $F$. The identity
of this corner need not be the identity of the whole group algebra;
adding its complement to both factors will give a scalar product equal
to $1$. This is an explicit instance of the finite-factor reduction
of Dykema and Juschenko~\cite[Theorem~2.2]{DykemaJuschenko2015}.

For the finite matrix block, let
\begin{equation}
\label{eq:finite-auxiliary-group}
B=(\mathbb Z/\ell)^m,
\qquad
F=B\rtimes S_m,
\end{equation}
where $S_m$ permutes the coordinates. For $1\le i\le m$ define
\begin{equation}
\label{eq:coordinate-idempotents}
p_i=|B|^{-1}\sum_{w\in B}\zeta^{-w_i}w\in K[F].
\end{equation}
The coordinate characters are distinct. The character-sum calculation
used in Equation~\eqref{eq:character-products}, now on $B$, shows that
the $p_i$ are pairwise orthogonal idempotents. Each is nonzero, since its
identity coefficient is $|B|^{-1}$. This denominator exists because
$|B|=\ell^m$ is odd. No inverse of $|S_m|$ is required.

Let $\sigma_i\in S_m$ exchange coordinates $1$ and $i$, taking
$\sigma_1=1$. Then $\sigma_i p_1\sigma_i^{-1}=p_i$. Put
\begin{equation}
\label{eq:finite-matrix-units}
p_{ij}=\sigma_i p_1\sigma_j^{-1}\qquad(1\le i,j\le m).
\end{equation}
These elements are nonzero, being left and right translates of $p_1$,
and they satisfy $p_{ii}=p_i$ and $p_i p_{ij}p_j=p_{ij}$.
For $j\ne k$ the latter identity and $p_jp_k=0$ give
$p_{ij}p_{kr}=0$. For $j=k$, direct cancellation gives
$p_{ij}p_{jr}=\sigma_i p_1^2\sigma_r^{-1}=p_{ir}$. Thus
\begin{equation}
\label{eq:matrix-unit-products}
p_{ij}p_{kr}=\delta_{jk}p_{ir}.
\end{equation}
The matrix units are linearly independent over $K$: multiplying a
linear relation on the left by $p_i$ and on the right by $p_j$ isolates
its coefficient of the nonzero element $p_{ij}$.

Take
\begin{equation}
\label{eq:output-group}
G=L\times F.
\end{equation}
We regard $K[L]$ and $K[F]$ as commuting subalgebras of $K[G]$, and put
\begin{equation}
\label{eq:block-embedding}
f_{\mathrm{blk}}=\sum_{i=1}^m p_{ii},
\qquad
\Phi:\Mat_m(K[L])\longrightarrow K[G],
\qquad
\Phi(Z)=\sum_{i,j=1}^m Z_{ij}p_{ij}.
\end{equation}
Equation~\eqref{eq:matrix-unit-products} proves that $\Phi$ is
multiplicative and that $\Phi(I_m)=f_{\mathrm{blk}}$.
It is also injective. Indeed, the group basis identifies
$K[L\times F]$ with $K[L]\otimes_K K[F]$. If $\Phi(Z)=0$, expand
each entry as $Z_{ij}=\sum_{g\in L}c_{ij,g}g$. Comparing the
$L$-basis coefficient of each $g$ gives
$\sum_{i,j}c_{ij,g}p_{ij}=0$. The linear independence just proved
forces every $c_{ij,g}$ to be zero.

The required two elements are
\begin{equation}
\label{eq:output-elements}
\boxed{\begin{aligned}
a_{\mathrm{out}}
 &=1-f_{\mathrm{blk}}+
   \sum_{i,j=1}^m
   \bigl((UJ_0)_{ij}+(1-e)\delta_{ij}\bigr)p_{ij},\\
b_{\mathrm{out}}
 &=1-f_{\mathrm{blk}}+
   \sum_{i,j=1}^m
   \bigl((CT)_{ij}+(1-e)\delta_{ij}\bigr)p_{ij}.
\end{aligned}}
\end{equation}
Equivalently, these are $1-f_{\mathrm{blk}}+\Phi(a_{\mathrm{mat}})$
and $1-f_{\mathrm{blk}}+\Phi(b_{\mathrm{mat}})$, respectively.
The matrix-unit identities imply
\[
f_{\mathrm{blk}}^2=f_{\mathrm{blk}},
\qquad
f_{\mathrm{blk}}\Phi(Z)=\Phi(Z)=\Phi(Z)f_{\mathrm{blk}}.
\]
Therefore $1-f_{\mathrm{blk}}$ annihilates every image of $\Phi$ on
both sides. It follows from Equation~\eqref{eq:square-forward-product}
that
\[
a_{\mathrm{out}}b_{\mathrm{out}}
=1-f_{\mathrm{blk}}+\Phi(I_m)=1,
\]
whereas
\begin{equation}
\label{eq:output-defect}
b_{\mathrm{out}}a_{\mathrm{out}}-1=\Phi(\Delta)\ne0
\end{equation}
by injectivity. Only the identity action of $f_{\mathrm{blk}}$ on the
image of $\Phi$ is needed for these computations.

\subsection{Finite presentation and the retained odd torsion}
\label{sec:finite-presentation}

All sums in Equation~\eqref{eq:output-elements} are finite, as are the
sums defining their entries. We now verify the remaining assertions about
the group $G$ by tracking its presentation through the construction.

Let $N=\langle\mathcal S_N\mid\mathcal R_N\rangle$ be the finite
presentation from Section~\ref{sec:criterion}, with the generators
$s_{f,c}$ and the local zero-sum relators of length at most three. With
$[x,y]=xyx^{-1}y^{-1}$, the direct product with the cyclic group has the
finite presentation
\[
H=\left\langle\mathcal S_N,z\ \middle|\;
\mathcal R_N,\ z^\ell,\ [z,s]\ (s\in\mathcal S_N)\right\rangle.
\]
Every element of $K_v$ and $K_*$ in
Equation~\eqref{eq:twisted-conjugation} has a representative given by
its local generator, or by the empty word for the identity. Using those
representatives, that equation adds $|V|$ stable letters and at most
$|V|\ell^3$ relators to the displayed presentation of $H$: its indices
are $v\in V$, $u\in K_v$, and $r\in\mathbb Z/\ell$. Hence $L$ is
finitely presented. The permutation action used above proves that the
base group embeds; its choices do not add generators or relators.

The group $F=(\mathbb Z/\ell)^m\rtimes S_m$ is finite, so its finite
multiplication table gives a finite presentation. A presentation of
$G=L\times F$ consists of finite presentations of the two factors on
disjoint generator sets, together with the commutators of every generator
of $L$ with every generator of $F$. This adds only finitely many relators,
and therefore $G$ is finitely presented. Finally, the embeddings already
proved give
\[
K_v\le N\le H\le L\le G.
\]
Since $K_v$ is the additive group of $\F_\ell^2$, its nonidentity
elements have order $\ell$. The output group therefore contains the
asserted odd-order torsion. This completes the proof of
Theorem~\ref{thm:criterion}.

The construction proves a failure of direct finiteness for a group
algebra by exhibiting its two elements. In particular, the last step
does not infer stable finiteness from direct finiteness for an arbitrary
ring: it uses the explicit finite group $F$ and the injective map
$\Phi$.
